\section{Introduction}

Thalassemia is an inherited hemoglobin disorder. Severe beta-thalassemia can
require lifelong red-cell transfusion and iron chelation. Current care has
credible standards, but curative and disease-modifying options remain unevenly
accessible because of donor availability, manufacturing infrastructure,
specialist centers, regulation, and cost \cite{farmakis2022tif}.

The practical cure-oriented endpoint for transfusion-dependent thalassemia is
durable transfusion independence with safe hemoglobin, controlled iron burden,
and acceptable long-term toxicity. Existing cure-like strategies point toward
hematopoietic stem cells: replace them through allogeneic transplantation,
modify them through lentiviral beta-globin gene addition, or edit them to
reactivate fetal hemoglobin \cite{locatelli2024exaCel,thompson2022beticel}.

Nakafa Lab treats these therapies as biological benchmarks and specialist
referral gates, not as repo-issued recommendations. The research question is
whether a lower-cost program can approach the same biology through fetal
hemoglobin induction, globin-chain balance, erythroid maturation, red-cell
survival, or iron-burden modification while remaining feasible in the Indonesian
access context.

\begin{figure}[htbp]
\centering
\begin{tikzpicture}[
  node distance=0.8cm and 1.1cm,
  box/.style={
    draw=blue!45!black,
    fill=blue!5,
    rounded corners=2pt,
    align=center,
    text width=0.25\textwidth,
    minimum height=1.1cm
  },
  arrow/.style={-{Stealth[length=2.4mm]}, thick, draw=blue!60!black}
]
\node[box] (sources) {Sources\\guidelines, PubMed, registries, regulators};
\node[box, right=of sources] (gates) {Evidence gates\\benchmark, candidate,
boundary};
\node[box, right=of gates] (assay) {Assay design\\HbF, globin balance,
hemolysis};
\node[box, below=of sources] (clinical) {Clinical packet\\subtype, immune, iron,
access};
\node[box, right=of clinical] (paper) {Public output\\paper, tables, repo};

\draw[arrow] (sources) -- (gates);
\draw[arrow] (gates) -- (assay);
\draw[arrow] (assay.south) |- (clinical.east);
\draw[arrow] (clinical) -- (paper);
\end{tikzpicture}
\caption{Nakafa Lab evidence-to-assay pipeline. The process is designed to make
rejection and uncertainty visible before patient-facing claims are made.}
\label{fig:evidence-pipeline}
\end{figure}

